\section{A rational Lie algebra of categorical values}\label{sec:category}

The upper bound on the full equation space is now established.  To
attain it, we construct a rational subspace that is closed under the
Ihara bracket and contains an element with nonzero coefficient of
$x^{n-1}y$ in every odd weight $n\geq3$.  This section proves the
closure property; the next constructs the nonzero values.
Only configurations with at most four labels will be needed.
We use the parenthesized-chord formalism of Bar-Natan
\cite[Sections~2--3]{BarNatan}, restricted to the finite arities and
operations specified below.
All completions in this section are by weight, and all tensor products of
completed spaces are completed by total weight.

We use the standard label sets $[k]=\{1,\ldots,k\}$ for
$1\leq k\leq4$; other label sets denote their relabeled copies.
For such a finite label set $I$, let
$\mathcal T(I)$ be the set of planar binary trees whose leaves are
bijectively labeled by $I$.  Thus the objects record both parentheses and
a linear ordering.  Put $\mathfrak t_I=0$ for $|I|=1$, and otherwise use
the infinitesimal braid algebra on $I$.  Over a field $K$ of
characteristic zero define a $K$-linear category $\mathcal C(I;K)$ by
\[
 \operatorname{Hom}(p,q)=\widehat U(\mathfrak t_I\otimes K)
 \qquad(p,q\in\mathcal T(I)).
\]
Composition is multiplication, with the last arrow on the left.  Write
$1_{pq}$ for the arrow with coefficient $1$, so that
$1_{qr}1_{pq}=1_{pr}$.  Each arrow space carries the coproduct determined
by
\[
 \Delta(t_{ij})=t_{ij}\otimes1+1\otimes t_{ij},
 \qquad \Delta(1_{pq})=1_{pq}\otimes1_{pq}.
\]
In particular, the unit coefficient arrow between different objects is
distinguished from an identity endomorphism.

Here are all the additional operations that we impose.  Relabeling is
allowed along every bijection of label sets.  For an outer tree on $I$
and fixed trees on disjoint nonempty sets $B_i$, graft the latter at its
leaves.  On arrows this operation is the algebra homomorphism
\begin{equation}\label{eq:category-block-sum}
 t_{ij}\longmapsto t_{B_iB_j}:=
 \sum_{a\in B_i,\ b\in B_j}t_{ab}.
\end{equation}
There is also insertion of a varying tree in one fixed leaf of a fixed
outer tree: on arrows this retains each inner chord $t_{ab}$ with its
labels.  Both operations are used whenever the resulting number of
leaves is at most four; single-leaf trees permit identity insertions.
These are functors preserving coproducts.  Indeed, the disjoint-chord
relations and the three-label relations imply the corresponding
relations between the block sums in
\eqref{eq:category-block-sum}, by summing first over labels in each
block.  Inner chords commute with outer block sums: for
$a,b\in B_i$ and $c\notin B_i$ the only possibly nonzero terms are
$[t_{ab},t_{ac}+t_{bc}]=0$.  These observations also show that successive
insertions agree with grafting the same trees in one step.

A \emph{compatible derivation of weight $n$} is a family $\delta$ of
$K$-linear maps on these arrow spaces, raising weight by $n$, such that
\[
 \delta(ba)=\delta(b)a+b\delta(a),\qquad
 \Delta\delta=(\delta\otimes\id+\id\otimes\delta)\Delta.
\]
It must commute with all the operations just specified and vanish on
all arrow spaces in arity two.  The arity-one derivation is necessarily
zero.  Let $\mathcal D_n(K)$ denote this space.  Commutators of such
families are again compatible derivations, of the sum of their weights.

Set
\[
 a=((12)3),\qquad b=(1(23)),\qquad
 \alpha=1_{ab},\qquad x=t_{12},\quad y=t_{23}.
\]
Evaluation on $\alpha$ will be understood as its coefficient in
$\widehat U(\mathfrak t_3\otimes K)$.

\begin{proposition}\label{prop:categorical-values}
For $n\geq2$, evaluation $\delta\mapsto\delta(\alpha)$ takes
$\mathcal D_n(\Q)$ into $W_n$.  Its images
\[
 V_n=\{\delta(\alpha):\delta\in\mathcal D_n(\Q)\},\qquad
 V=\bigoplus_{n\geq2}V_n,
\]
form a rational Lie subalgebra under the Ihara bracket \eqref{eq:ihara}.
\end{proposition}

\begin{proof}
Write $\psi=\delta(\alpha)$.  The coderivation rule at the group-like
arrow $\alpha$ says
\[
 \Delta\psi=\psi\otimes1+1\otimes\psi.
\]
Primitives in a characteristic-zero enveloping algebra are its Lie
algebra: by the PBW filtration, a primitive of filtration degree
$d>1$ would give a primitive homogeneous polynomial of degree $d$ in a
symmetric algebra, whereas its coproduct has a nonzero component of
bidegree $(1,d-1)$.  Induction on filtration degree proves the assertion.
Now
\begin{equation}\label{eq:category-three-split}
 \mathfrak t_3=\Lie_K\langle x,y\rangle\oplus K T,
 \qquad T=t_{12}+t_{13}+t_{23},
\end{equation}
with $T$ central: replacing $t_{13}$ by $T-x-y$ transforms exactly the
three defining relations into the centrality of $T$.  Thus a primitive
of weight $n\geq2$ is a Lie polynomial $\psi(x,y)$.

Reversing the two children at any fork is insertion, followed if
necessary by an inner insertion, of a unit arrow in arity two.  Its
$\delta$-value is zero.  The values of unit coefficient arrows add under
composition, and inverse unit arrows have opposite values.  Relabel
$\alpha$ by $1\leftrightarrow3$.  Its source $((32)1)$ is carried to $b$
by fork reversals, and its target $(3(21))$ is carried to $a$ in the same
way.  Its value is therefore both $\psi(y,x)$ and $-\psi(x,y)$, giving
\begin{equation}\label{eq:category-antisymmetry}
 \psi(x,y)+\psi(y,x)=0.
\end{equation}
For the cyclic relation, use the successive trees
$((12)3),((23)1),((31)2)$.  Reassociation on the ordered triples
$(1,2,3),(2,3,1),(3,1,2)$, each followed by a root-fork reversal,
forms a cycle of unit arrows.  Consequently
\[
 \psi(t_{12},t_{23})+\psi(t_{23},t_{31})+
 \psi(t_{31},t_{12})=0.
\]
An occurrence of the central $T$ in a Lie monomial of degree greater
than one contributes zero.  Substitution of $t_{31}=T-x-y$ gives
\begin{equation}\label{eq:category-cyclic}
 \psi(x,y)+\psi(y,-x-y)+\psi(-x-y,x)=0.
\end{equation}

Figure~\ref{fig:reassociation} shows the five four-leaf trees and
the two paths of reassociation arrows:
\[
\begin{gathered}
 A=(((12)3)4),\quad B=((1(23))4),\quad C=(1((23)4)),\\
 D=((12)(34)),\qquad E=(1(2(34))),\\
 A\longrightarrow D\longrightarrow E,
 \qquad A\longrightarrow B\longrightarrow C\longrightarrow E.
\end{gathered}
\]
The two edge values on the first path are
$\psi(t_{13}+t_{23},t_{34})$ and
$\psi(t_{12},t_{23}+t_{24})$.  Those on the second are
$\psi(t_{12},t_{23})$,
$\psi(t_{12}+t_{13},t_{24}+t_{34})$, and
$\psi(t_{23},t_{34})$.
Both composites are $1_{AE}$, hence
\begin{equation}\label{eq:category-pentagon}
\begin{split}
 \psi(t_{12},t_{23}+t_{24})+
 \psi(t_{13}+t_{23},t_{34})
 ={}&\psi(t_{23},t_{34})\\
 &+\psi(t_{12}+t_{13},t_{24}+t_{34})
 +\psi(t_{12},t_{23}).
\end{split}
\end{equation}
This proves membership in $W_n$ using only the stated operations.

\begin{figure}[!t]
  \centering
  \begin{tikzpicture}[
    x=1cm,y=1cm,
    vertex/.style={minimum width=2.7cm,minimum height=1.9cm,inner sep=0pt},
    tree/.style={line width=.65pt,line cap=round,line join=round},
    reassociate/.style={-{Stealth[length=2mm]},line width=.7pt},
    edge value/.style={fill=white,inner sep=3pt,font=\small},
    leaf/.style={font=\scriptsize,inner sep=1pt},
    tree caption/.style={font=\small,inner sep=1pt}
  ]
    \node[vertex] (A) at (-4.2,0) {};
    \node[vertex] (D) at (0,2.4) {};
    \node[vertex] (E) at (4.2,0) {};
    \node[vertex] (B) at (-2.7,-3) {};
    \node[vertex] (C) at (2.7,-3) {};

    \draw[reassociate] (A) -- node[edge value,above left] {$p_1$} (D);
    \draw[reassociate] (D) -- node[edge value,above right] {$p_2$} (E);
    \draw[reassociate] (A) -- node[edge value,left] {$p_3$} (B);
    \draw[reassociate] (B) -- node[edge value,above] {$p_4$} (C);
    \draw[reassociate] (C) -- node[edge value,right] {$p_5$} (E);

    \begin{scope}[shift={(A)}]
      \draw[tree] (0,.75) -- (-.3,.4) -- (-.6,.05) -- (-.9,-.3)
        (-.6,.05) -- (-.3,-.3) (-.3,.4) -- (.3,-.3)
        (0,.75) -- (.9,-.3);
      \foreach \x/\i in {-.9/1,-.3/2,.3/3,.9/4}
        \node[leaf,below] at (\x,-.3) {$\i$};
      \node[tree caption] at (0,-.78) {$\mathsf{A}:\ (((12)3)4)$};
    \end{scope}
    \begin{scope}[shift={(D)}]
      \draw[tree] (0,.75) -- (-.6,.2) -- (-.9,-.3)
        (-.6,.2) -- (-.3,-.3)
        (0,.75) -- (.6,.2) -- (.3,-.3)
        (.6,.2) -- (.9,-.3);
      \foreach \x/\i in {-.9/1,-.3/2,.3/3,.9/4}
        \node[leaf,below] at (\x,-.3) {$\i$};
      \node[tree caption] at (0,-.78) {$\mathsf{D}:\ ((12)(34))$};
    \end{scope}
    \begin{scope}[shift={(E)}]
      \draw[tree] (0,.75) -- (-.9,-.3)
        (0,.75) -- (.3,.4) -- (-.3,-.3)
        (.3,.4) -- (.6,.05) -- (.3,-.3)
        (.6,.05) -- (.9,-.3);
      \foreach \x/\i in {-.9/1,-.3/2,.3/3,.9/4}
        \node[leaf,below] at (\x,-.3) {$\i$};
      \node[tree caption] at (0,-.78) {$\mathsf{E}:\ (1(2(34)))$};
    \end{scope}
    \begin{scope}[shift={(B)}]
      \draw[tree] (0,.75) -- (-.3,.4) -- (-.9,-.3)
        (-.3,.4) -- (0,.05) -- (-.3,-.3)
        (0,.05) -- (.3,-.3) (0,.75) -- (.9,-.3);
      \foreach \x/\i in {-.9/1,-.3/2,.3/3,.9/4}
        \node[leaf,below] at (\x,-.3) {$\i$};
      \node[tree caption] at (0,-.78) {$\mathsf{B}:\ ((1(23))4)$};
    \end{scope}
    \begin{scope}[shift={(C)}]
      \draw[tree] (0,.75) -- (-.9,-.3)
        (0,.75) -- (.3,.4) -- (0,.05) -- (-.3,-.3)
        (0,.05) -- (.3,-.3) (.3,.4) -- (.9,-.3);
      \foreach \x/\i in {-.9/1,-.3/2,.3/3,.9/4}
        \node[leaf,below] at (\x,-.3) {$\i$};
      \node[tree caption] at (0,-.78) {$\mathsf{C}:\ (1((23)4))$};
    \end{scope}
    \node[align=center,font=\small] at (0,-.55)
      {$p_1+p_2=p_3+p_4+p_5$};
  \end{tikzpicture}

  \smallskip
  {\small
  \begin{tabular}{@{}cl@{\qquad}cl@{}}
    \toprule
    Label & Infinitesimal value & Label & Infinitesimal value \\
    \midrule
    $p_1$ & $\psi(t_{13}+t_{23},t_{34})$
      & $p_3$ & $\psi(t_{12},t_{23})$ \\
    $p_2$ & $\psi(t_{12},t_{23}+t_{24})$
      & $p_4$ & $\psi(t_{12}+t_{13},t_{24}+t_{34})$ \\
    & & $p_5$ & $\psi(t_{23},t_{34})$ \\
    \bottomrule
  \end{tabular}}
  \caption{The five parenthesizations of four ordered leaves. Each arrow
  reassociates $((IJ)K)$ to $(I(JK))$, with the appropriate fixed subtrees
  inserted. The labels $p_1,\ldots,p_5$ denote the infinitesimal values
  in the table, not the arrows themselves.
  The two paths $\mathsf{A}\to\mathsf{D}\to\mathsf{E}$ and
  $\mathsf{A}\to\mathsf{B}\to\mathsf{C}\to\mathsf{E}$ have the same unit
  arrow as composite. Their infinitesimal values add, yielding the pentagon
  identity.}
  \label{fig:reassociation}
\end{figure}


It remains to check the bracket, including its sign.  Let $d_a,d_b$
be the loop derivations at $a,b$.  Inner insertion gives $d_a(x)=0$
and $d_b(y)=0$.  Differentiating the equality of arrows
$\alpha y_a=y_b\alpha$ gives
\[
 \psi y+d_a(y)=y\psi,
 \qquad d_a(y)=[y,\psi].
\]
Outer arity-two insertion also kills $t_{13}+t_{23}$ at $a$, so
$d_a(T)=0$.  In the decomposition
\eqref{eq:category-three-split} its restriction is exactly $D_\psi$.
An arrow from $a$ to $b$ with coefficient $h$ equals $\alpha h_a$;
therefore
\begin{equation}\label{eq:category-arrow-derivative}
 \delta_{ab}(h)=\psi h+D_\psi(h)
 \qquad(h\in K\langle\!\langle x,y\rangle\!\rangle).
\end{equation}
If another derivation $\varepsilon$ has value $\phi$, their commutator
has value
\[
 [\delta,\varepsilon](\alpha)
 =D_\psi(\phi)-D_\phi(\psi)+\psi\phi-\phi\psi
 =\{\psi,\phi\}.
\]
Evaluation is thus a Lie homomorphism onto its image.  No injectivity
of evaluation is required.
\end{proof}

\begin{lemma}[Rational finite data]\label{lem:rational-descent}
For every $n>0$ and every characteristic-zero extension $K/\Q$,
\[
 \mathcal D_n(K)=\mathcal D_n(\Q)\otimes_{\Q}K.
\]
For $n\geq2$, the image of evaluation is $V_n\otimes_{\Q}K$.
In particular, if the coefficient of a fixed associative word is
nonzero on a complex categorical value of weight $n$, it is nonzero
on some rational categorical value of that weight.
\end{lemma}

\begin{proof}
At each arity choose a base object $o$.  For a composition derivation
put $h_p=\delta(1_{op})$, with $h_o=0$, and let $d$ be its derivation on
loops at $o$.  Since every arrow has the factorization
$1_{oq}\,u_o\,1_{po}$, the product and inverse rules force
\begin{equation}\label{eq:category-finite-data}
 \delta_{pq}(u)=h_qu+d(u)-uh_p.
\end{equation}
Conversely, an algebra derivation $d$ and elements $h_p$ define a
composition derivation by this formula.  In weight $n$ the data are
\[
 h_p\in U(\mathfrak t_I)_n,
 \qquad d(t_{ij})\in U(\mathfrak t_I)_{n+1}.
\]
There are finitely many objects and generators, and every displayed
space is finite dimensional over $\Q$.

The conditions that $d$ preserve the defining braid relations are
rational linear equations in these data.  The coderivation conditions
are equivalent to primitivity of all $h_p$ and $d(t_{ij})$: necessity
follows on the reference unit arrows and generator loops, and
sufficiency follows by the product rule and
\eqref{eq:category-finite-data}.  These are again finitely many rational
linear equations, since coproducts in the indicated weights take
values in finite sums of finite-dimensional spaces.

Finally, there are only finitely many label permutations, object
graftings, and inner placements through arity four.  For each associated
functor $\rho$, the condition $\delta\rho=\rho\delta$ need only be
checked on the reference unit arrows and generator loops: all other
arrows follow by composition, the product rule, and continuity.
These checks, and the requirement of zero arity-two action, are
rational linear equations.  Thus $\mathcal D_n$ is the kernel of a
linear map between finite-dimensional rational spaces.  Kernels and
images of such maps commute with scalar extension.  Evaluation and
coefficient extraction are rational linear maps, proving the last
assertion as well.
\end{proof}
